\documentclass[]{manual}

\usepackage{lenses}

\title{Thick Lenses}
\date{\today}
\course{Physics 325}

\begin{document}
    \maketitle
    \thispagestyle{firstpage}

    %%%

    \section*{Motivation}

    Although thin lenses are far more common, thick lenses have some situationally desirable attributes.
    For example, they can provide much higher optical power ($P = \frac{1}{f}$) with less delicate construction than a multi-thin-lens optical system.
    For example, the thick lenses we will work with in this lab are World War 2 bomb sights.
    A single thick lens was (apparently) much easier to design around than what would likely be a complicated arrangement of thin lenses.
    People with particularly poor eyesight typically wear glasses with very thick lenses - again, it would be much harder to build glasses that held a system of thin lenses with the same optical power.

    \section*{Background}
    
    \begin{align}
        \frac{1}{F} &= \frac{1}{f_1} + \frac{1}{f_2} - \frac{a}{f_1 \, f_2} \\
        b_1 &= \frac{a \, F}{f_2} \\
        b_2 &= \frac{a \, F}{f_1} \\
    \end{align}
    
    \begin{figure}
        \centering
        
        \begin{tikzpicture}
            \draw[dashed] (-7, 0) -- (8, 0) node[left, at start] {$\mathrm{OA}$};
            
            \clensLabelled{-1}{40}{4}{0.05}{}
            \clensLabelled{1}{40}{4}{0.05}{}
            
            \draw[fill] (-3.5, 0) circle (2pt) node[above right] (PFO) {$\mathrm{PF}_1$};
            \draw[fill] (3.5, 0) circle (2pt) node[above right] (PFI) {$\mathrm{PF}_2$};
            
            \draw[<->, dashed] (-1, -4.5) -- ++(2, 0) node[midway, below] {$a$};
            
            \draw[dashed, thick] (-.5, -3.5) -- ++(0, 9) node[above] {$\mathrm{PP}_1$};
            \draw[dashed, thick] (.5, -3.5) -- ++(0, 9) node[above] {$\mathrm{PP}_2$};
            
            \draw[<->|, dashed] (0.5, -.7) -- (4, -.7) node[midway, below] (pfi) {$F$};
            \draw[<->|, dashed] (-0.5, -.7) -- (-4, -.7) node[midway, below] (pfo) {$F$};
            
            \object{-6}{2}
    
            \draw[neararrow={latex}] (-6, 1.8) -- (-0.5, 1.8);
            \draw[midarrow={latex}, dotted] (-0.5, 1.8) -- ++(1, 0);
            \draw[neararrow={latex}, shorten >= -5cm] (0.5, 1.8) -- (3.5, 0);
            
            \pgfmathsetmacro{\y}{{1.8*3/2.5}}
            \draw[veryneararrow={latex}] (-6, 1.8) -- (-.5, -\y);
            \draw[midarrow={latex}, dotted] (-.5, -\y) -- ++(1, 0);
            \draw[neararrow={latex}] (0.5, -\y) -- ++(7.5, 0);
            
            \draw[-{squareTip[scale=3]}] ({0.5+1/(1/3-1/5.5)}, 0) -- ++(0, -2.35) node[below] {Image};
        \end{tikzpicture}
        
        \caption{
            The effective thick lens optical system, made from two thin lenses.
        }
        
        \label{fig:effective-thick-lens}
    \end{figure}


    \section*{Equipment}


    \FloatBarrier
    \newpage

    \section*{Procedure}
    
    \begin{figure}
        \centering
        
        \begin{tikzpicture}
            \draw[dashed] (-6, 0) -- (4, 0) node[left, at start] {$\mathrm{OA}$};
        
            \pgfmathsetmacro{\h}{2}
            \draw[ultra thick] (-2, \h) -- ++(0, -2*\h)
            -- ++(2, 0)
            .. controls (1, -0.5) and (1, 0.5) .. ++(0, 2*\h)
            node[pos=.125, inner sep=0] (a) {}
            node[pos=.25, inner sep=0] (b) {}
            node[pos=.375, inner sep=0] (c) {}
            node[pos=.5, inner sep=0] (d) {}
            node[pos=.625, inner sep=0] (e) {}
            node[pos=.75, inner sep=0] (f) {}
            node[pos=.875, inner sep=0] (g) {}
            -- cycle node[above, midway] {Thick Lens};
            
            \draw[ultra thick, fill] (4, 2) -- ++(0, -4) -- ++(0.1, 0) -- ++(0, 4) -- cycle node[above, midway] {Screen};
            
            \foreach \y in {-1.5,-1,...,1.5} {
                \draw[midarrow={latex}] (-5, \y) -- (-2, \y);
            }
            \draw[midarrow={latex}] (a) -- (4, 0);
            \draw[midarrow={latex}] (b) -- (4, 0);
            \draw[midarrow={latex}] (c) -- (4, 0);
            \draw[midarrow={latex}] (d) -- (4, 0);
            \draw[midarrow={latex}] (e) -- (4, 0);
            \draw[midarrow={latex}] (f) -- (4, 0);
            \draw[midarrow={latex}] (g) -- (4, 0);
            
            \node[above] (object) at (-5, 1.6) {Rays from Object at $\infty$};
        \end{tikzpicture}
        
        \caption{
            Imaging an object at infinity using the thick lens.
        }
        
        \label{fig:object-at-infinity}
    \end{figure}


    \FloatBarrier
    \newpage

    \section*{Prelab Questions}
    \setcounter{subsection}{0}

    \subsection{Question 1 \hfill (1 pt.)}

    \vspace{\stretch{1}}

    \subsection{Question 2 \hfill (1 pt.)}

    \vspace{\stretch{1}}

\end{document}
